% ===========================================================================
% Physical AI Systems, LaTeX preamble
% ===========================================================================
%
% Engine:  LuaLaTeX
% Class:   KOMA-Script scrbook (set in _quarto.yml)
% Loaded:  first of the include-in-header files. tex/cover.tex follows and
%          depends on the palette defined here.
%
% RULES FOR EDITING THIS FILE
%
%   1. Every setting is declared exactly once. To change a value, change it
%      where it lives. Do not append an override at the bottom. This file
%      previously carried two \setlist declarations, two \arraystretch values,
%      two \emergencystretch values, and fvextra loaded twice, all because
%      changes were bolted on rather than integrated.
%   2. Sections run in dependency order: packages, palette, page, type,
%      floats, code, notes, links, utilities. Put new settings in the section
%      they belong to.
%   3. Anything that draws something belongs in cover.tex or in a figure
%      source, not here.
%   4. Comment the intent, not the syntax. A future reader needs to know why a
%      value is what it is, not that \setlength sets a length.
%
% Page geometry, font size, and leading live in _quarto.yml, because Quarto
% owns the geometry package options. This file assumes them, never sets them.
%
% ===========================================================================


% ---------------------------------------------------------------------------
% 1. Packages
% ---------------------------------------------------------------------------

\usepackage{microtype}          % Optical margins and glyph expansion.
\usepackage{float}              % Strict float placement [H].
\usepackage{morefloats}         % Expand LaTeX float registers.
\usepackage{placeins}           % Prevent floats from overshooting section/chapter boundaries.
\usepackage{needspace}          % Dynamic space reservation to prevent awkward table breaks.
\usepackage{scrlayer-scrpage}   % KOMA page styles, for running heads.
\usepackage{etoolbox}           % Patching KOMA and Quarto internals.
\usepackage{array}              % Column specifiers used by pandoc tables.
\usepackage{booktabs}           % \toprule and friends in pandoc tables.
\usepackage{enumitem}           % List spacing control.
\usepackage{framed}             % Backing for the Shaded code environment.
\usepackage{fancyvrb}           % Verbatim engine that pandoc emits into.
\usepackage{fvextra}            % Line breaking inside listings.
\usepackage{ragged2e}            % Balanced ragged-right typesetting with hyphenation.
\usepackage{marginnote}         % Margin column content.
\usepackage{environ}            % \NewEnviron, used by the fullwidth utility.
\usepackage{lettrine}           % Drop caps for chapter openers.
\usepackage{colortbl}           % Alternating row tints and table styling.
\usepackage{tikz}               % Cover artwork.
\usepackage{tikzpagenodes}      % Margin and page positioning for chapter numbers.
\usetikzlibrary{arrows.meta,positioning,shapes.geometric,fit,backgrounds,calc}


% ---------------------------------------------------------------------------
% 2. Palette
%
% Harvard crimson and ETH Zurich blue, plus a semantic layer. Prose, callouts,
% and figures reference the semantic names, so a palette change happens in one
% place. The raw names stay because the TikZ figure sources use them directly.
% ---------------------------------------------------------------------------

\definecolor{harvardcrimson}{HTML}{A51C30}
\definecolor{ethdarkblue}{HTML}{1F407A}
\definecolor{ethblue}{HTML}{215CAF}
\definecolor{ethpetrol}{HTML}{007A87}
\definecolor{ethbronze}{HTML}{B87333}
\definecolor{ethpurple}{HTML}{5B4B8A}

% Semantic layer. These encode the book's argument: what is structural, what
% is measured, what is only proposed, what has consequences, who has authority.
\colorlet{structural}{ethdarkblue}
\colorlet{measured}{ethpetrol}
\colorlet{proposal}{ethbronze}
\colorlet{consequence}{harvardcrimson}
\colorlet{authority}{ethpurple}

% Ink levels for body, secondary, and hairline material.
\definecolor{ink}{HTML}{1A202C}
\definecolor{softink}{HTML}{4A5568}
\definecolor{paleink}{HTML}{A0AEC0}

% Surfaces. shadecolor is named by the framed/Shaded contract, not by us.
\definecolor{cardbg}{HTML}{F8FAFC}
\definecolor{shadecolor}{HTML}{F4F6F9}

% Quarto's standard callouts, harmonized to the semantic layer.
\definecolor{quarto-callout-note-color}{HTML}{1F407A}
\definecolor{quarto-callout-note-color-frame}{HTML}{1F407A}
\definecolor{quarto-callout-tip-color}{HTML}{007A87}
\definecolor{quarto-callout-tip-color-frame}{HTML}{007A87}
\definecolor{quarto-callout-warning-color}{HTML}{A51C30}
\definecolor{quarto-callout-warning-color-frame}{HTML}{A51C30}
\definecolor{quarto-callout-important-color}{HTML}{5B4B8A}
\definecolor{quarto-callout-important-color-frame}{HTML}{5B4B8A}
\definecolor{quarto-callout-caution-color}{HTML}{A51C30}
\definecolor{quarto-callout-caution-color-frame}{HTML}{A51C30}


% ---------------------------------------------------------------------------
% 3. Page furniture (Running headers & footers)
%
% Standard two-sided academic textbook layout (Cambridge / MIT Press):
%   Even (verso / left) pages:  [Page #] ........................ [Chapter Title]
%   Odd (recto / right) pages:  [Section Title] ........................ [Page #]
% Subtitles and linebreaks are strictly stripped from headers.
% ---------------------------------------------------------------------------

\clearpairofpagestyles

\KOMAoptions{headwidth=textwithmarginpar,footwidth=textwithmarginpar}
\KOMAoptions{headsepline=0.4pt}
\KOMAoptions{markcase=noupper}
\addtokomafont{headsepline}{\color{accentcolor}}

\renewcommand*{\chaptermarkformat}{}
\renewcommand*{\sectionmarkformat}{}

\setkomafont{pagehead}{\normalfont\rmfamily\fontsize{8.5pt}{10.5pt}\selectfont\color{accentcolor}}
\setkomafont{pagenumber}{\normalfont\rmfamily\fontsize{8.5pt}{10.5pt}\selectfont\color{accentcolor}}

% Verso (even / left) header: Left is Folio, Right is Section mark (matching MLSysBook)
\lehead{\pagemark}
\rehead{\small\rightmark}

% Recto (odd / right) header: Left is Chapter mark, Right is Folio (matching MLSysBook)
\lohead{\small\leftmark}
\rohead{\pagemark}

% Chapter opening pages (plain style): centered folio at bottom in accent color
\cfoot[\pagemark]{}
\setkomafont{pagefoot}{\normalfont\rmfamily\fontsize{8.5pt}{10.5pt}\selectfont\color{accentcolor}}

\automark[section]{chapter}



% ---------------------------------------------------------------------------
% 4. Typography
%
% Serif body at a 4.85in measure, with matching classical serif headings
% matching MLSysBook Volumes 1 and 2.
%
% Paragraphs are indented rather than separated by vertical space. Vertical
% space marks structure here, not every paragraph break.
% ---------------------------------------------------------------------------

\KOMAoptions{fontsize=9.5pt}

\frenchspacing
\microtypesetup{protrusion=true, expansion=true, final}

\setlength{\parskip}{0pt plus 0.6pt}
\setlength{\parindent}{1.2em}

% Justification. Long technical compounds at this measure produce ladders of
% hyphens. Loosen the badness tolerance rather than the word spacing, so the
% grey of the page stays even.
\hyphenpenalty=180
\tolerance=1400
\setlength{\emergencystretch}{2.2em}

% Widow and orphan control (Cambridge / MIT Press standard)
\clubpenalty=10000
\widowpenalty=10000
\displaywidowpenalty=10000
\raggedbottom

% Pandoc emits \linebreak inside margin material where it is invalid.
\let\origlinebreak\linebreak
\renewcommand{\linebreak}[1][4]{\ifhmode\origlinebreak[#1]\fi}

% --- Headings (Matching MLSysBook Volumes 1 and 2) ------------------------

\setkomafont{disposition}{\normalfont\rmfamily\color{accentcolor}}
\setkomafont{chapter}{\normalfont\rmfamily\bfseries\fontsize{20pt}{24pt}\selectfont\color{accentcolor}}
\setkomafont{section}{\normalfont\rmfamily\bfseries\fontsize{13pt}{16pt}\selectfont\color{accentcolor}}
\setkomafont{subsection}{\normalfont\rmfamily\bfseries\fontsize{11pt}{14pt}\selectfont\color{accentcolor}}
\setkomafont{subsubsection}{\normalfont\rmfamily\bfseries\fontsize{10pt}{13pt}\selectfont\color{accentcolor}}
\setkomafont{paragraph}{\normalfont\rmfamily\bfseries\fontsize{9.5pt}{12pt}\selectfont\color{accentcolor}}

% Drop caps configuration for chapter openings (Matching MLSysBook)
\renewcommand{\LettrineFontHook}{\rmfamily\bfseries\color{accentcolor}}
\renewcommand{\LettrineTextFont}{\normalfont\scshape}
\setlength{\DefaultFindent}{2pt}
\setlength{\DefaultNindent}{0pt}

% Display math spacing
\setlength{\abovedisplayskip}{6pt plus 2pt minus 2pt}
\setlength{\belowdisplayskip}{6pt plus 2pt minus 2pt}
\setlength{\abovedisplayshortskip}{0pt plus 2pt}
\setlength{\belowdisplayshortskip}{4pt plus 2pt minus 2pt}

% Chapter subtitle command for chapter openers (appears before the horizontal rule)
\DeclareRobustCommand{\chaptersubtitle}[1]{%
  \texorpdfstring{%
    \newline{\normalfont\fontsize{11pt}{14pt}\selectfont\color{ethdarkblue}\textit{#1}}%
  }{}%
}

% --- Chapter & Section Formatting (Matching Volume 1 & Volume 2) ---------
\usepackage[explicit]{titlesec}

% Giant chapter-number font in margin
\newcommand{\mlsysChapNumFont}{\normalfont\rmfamily\fontsize{70pt}{70pt}\selectfont\bfseries}

% Numbered chapters (Large serif title with giant margin number)
\titleformat{\chapter}[block]
  {\normalfont\rmfamily\bfseries\color{accentcolor}}
  {}
  {0pt}
  {%
    \begin{tikzpicture}[remember picture, overlay]
      \node[anchor=north east,
            inner sep=0pt,
            font={\mlsysChapNumFont},
            text=accentcolor]
        at ([xshift=-40pt]current page.north east |- current page text area.north)
        {\thechapter};
    \end{tikzpicture}%
    {\fontsize{34pt}{38pt}\selectfont #1}%
  }
  []

% Unnumbered chapters (Preface, Notation, Syllabus): clean large serif title without horizontal bars
\titleformat{name=\chapter,numberless}
  {\normalfont\rmfamily\huge\bfseries\color{accentcolor}}
  {}
  {0pt}
  {\Huge #1}
  []

\titlespacing*{\chapter}{0pt}{0pt}{24pt}

% Section (Large, bold serif, accentcolor)
\titleformat{\section}
  {\normalfont\rmfamily\Large\bfseries\color{accentcolor}\raggedright}
  {\thesection}
  {0.5em}
  {#1}
\titlespacing*{\section}{0pc}{12pt plus 3pt minus 3pt}{4pt plus 1pt minus 1pt}[0pc]

% Subsection (large, bold serif, accentcolor)
\titleformat{\subsection}
  {\normalfont\rmfamily\large\bfseries\color{accentcolor}\raggedright}
  {\thesubsection}
  {0.5em}
  {#1}
\titlespacing*{\subsection}{0pc}{10pt plus 3pt minus 3pt}{4pt plus 1pt minus 1pt}[0pc]

% Subsubsection (normal size, bold serif, accentcolor)
\titleformat{\subsubsection}
  {\normalfont\rmfamily\normalsize\bfseries\color{accentcolor}\raggedright}
  {\thesubsubsection}
  {0.5em}
  {#1}
\titlespacing*{\subsubsection}{0pc}{9pt plus 3pt minus 3pt}{3pt plus 1pt minus 1pt}[0pc]

% --- Lists ----------------------------------------------------------------

\setlist[itemize]{leftmargin=1.1em,topsep=2pt,itemsep=1pt,parsep=0pt}
\setlist[enumerate]{leftmargin=1.3em,topsep=2pt,itemsep=1pt,parsep=0pt}

% --- Table of Contents & Structure Dividers (MLSysBook Modern Layout) ------

\makeatletter
% State flags for TOC divisions
\newif\ifmlsysinbackmatter
\mlsysinbackmatterfalse
\newif\ifsynthesiswritten
\synthesiswrittenfalse
\newif\if@firstnumbered
\@firstnumberedtrue
\newcounter{mlsysappendixcount}

% Pure TeX helper to distinguish Chapter numbers (1..17) from Appendix letters (A..G)
\newcommand{\format@toc@number}[1]{%
  \if\relax\detokenize\expandafter{\romannumeral-0#1}\relax
    Chapter~#1%
  \else
    Appendix~#1%
  \fi
}

% Start frontmatter with Roman numerals and suppress subsection noise in frontmatter
\AtBeginDocument{%
  \pagenumbering{roman}%
  \renewcommand{\thepage}{\roman{page}}%
  \renewcommand*\contentsname{{\normalfont\bfseries\fontsize{18pt}{22pt}\selectfont\color{consequence}Contents}}%
  \addtocontents{toc}{\protect\setcounter{tocdepth}{0}}%
  % Hook addcontentsline AFTER hyperref initialization
  \let\orig@addcontentsline\addcontentsline
  \renewcommand{\addcontentsline}[3]{%
    \Ifstr{#1}{toc}{%
      \Ifstr{#2}{part}{%
        \ifmlsysinbackmatter
          % Suppress Appendices part divider; Back Matter division handles it
        \else
          \orig@addcontentsline{#1}{#2}{#3}%
        \fi
      }{%
        \Ifstr{#2}{chapter}{%
          \ifmlsysinbackmatter
            \Ifstr{#3}{References}{%
              \orig@addcontentsline{#1}{#2}{#3}%
            }{%
              \Ifstr{#3}{Index}{%
                \orig@addcontentsline{#1}{#2}{#3}%
              }{%
                \stepcounter{mlsysappendixcount}%
                \orig@addcontentsline{toc}{chapter}{\protect\numberline{\Alph{mlsysappendixcount}}#3}%
              }%
            }%
          \else
            \orig@addcontentsline{#1}{#2}{#3}%
          \fi
        }{%
          \orig@addcontentsline{#1}{#2}{#3}%
        }%
      }%
    }{%
      \orig@addcontentsline{#1}{#2}{#3}%
    }%
  }%
}

\let\orig@appendix\appendix
\renewcommand{\appendix}{%
  \orig@appendix
  \global\mlsysinbackmattertrue
  \addtocontents{toc}{\protect\tocdivisionentry{Back Matter}}%
  \addtocontents{toc}{\protect\setcounter{tocdepth}{0}}%
}

% Patch \mainmatter so Quarto does not prematurely switch to Arabic before Chapter 1
\let\orig@mainmatter\mainmatter
\renewcommand{\mainmatter}{%
  % Defer pagenumbering{arabic} until Chapter 1 is reached
}

\newif\ifiscontentstoc
\iscontentstocfalse

% Patch \tableofcontents to guarantee Crimson/Accent "Contents" heading
\let\orig@tableofcontents\tableofcontents
\renewcommand{\tableofcontents}{%
  \begingroup
    \iscontentstoctrue
    \DeclareRobustCommand{\chaptersubtitle}[1]{}%
    \renewcommand*\contentsname{{\normalfont\huge\bfseries\color{accentcolor}Contents}}%
    \orig@tableofcontents
  \endgroup
}

% Major division banners in TOC (e.g. "Main Text", "Synthesis", "Back Matter")
\DeclareRobustCommand{\tocdivisionentry}[1]{%
  \par\addvspace{18pt}%
  \noindent\hfil{\normalfont\sffamily\bfseries\fontsize{10.5pt}{13.5pt}\selectfont\color{ink}#1}\hfil\par%
  \vspace{3pt}%
  \noindent\hfil{\color{paleink!80!softink}\rule{0.45\textwidth}{0.45pt}}\hfil\par%
  \addvspace{8pt}%
}

% Part banners in TOC (e.g. "Part I: The Machine")
\DeclareRobustCommand{\tocpartentry}[1]{%
  \par\addvspace{16pt}%
  \noindent\hfil{\normalfont\sffamily\bfseries\fontsize{10pt}{13pt}\selectfont\color{consequence}#1}\hfil\par%
  \addvspace{7pt}%
}

% Hook into \addparttocentry to emit clean centered Crimson Part entries and division lines
\renewcommand*{\addparttocentry}[2]{%
  \Ifstr{#1}{}{%
    \Ifstr{#2}{Appendices}{}{%
      \addtocontents{toc}{\protect\tocdivisionentry{#2}}%
    }%
  }{%
    \Ifstr{#2}{}{%
      \addtocontents{toc}{\protect\tocpartentry{Part~#1}}%
    }{%
      \addtocontents{toc}{\protect\tocpartentry{#2}}%
    }%
  }%
}

% Suppress raw Appendices part entry and format part lines cleanly in TOC
\renewcommand*{\l@part}[2]{%
  \in@{Part}{#1}%
  \ifin@
    \par\addvspace{16pt}%
    \noindent\hfil{\normalfont\sffamily\bfseries\fontsize{10pt}{13pt}\selectfont\color{consequence}%
      \begingroup
        \renewcommand{\numberline}[1]{}%
        #1%
      \endgroup
    }\hfil\par%
    \addvspace{7pt}%
  \fi
}

% Intercept \chapter command to cleanly handle frontmatter vs mainmatter transitions
\let\old@chapter\chapter
\renewcommand{\chapter}{%
  \@ifstar{\unnumbered@chapter}{\numbered@chapter}%
}

\newcommand{\cleanmark}[1]{%
  \begingroup
    \def\chaptersubtitle##1{}%
    \def\newline{}%
    \def\texorpdfstring##1##2{##2}%
    \protected@xdef\@clean@title{#1}%
  \endgroup
  \markboth{\@clean@title}{}%
}

\newcommand{\unnumbered@chapter}[1]{%
  \cleardoublepage
  \old@chapter*{#1}%
  \cleanmark{#1}%
}

\newcommand{\numbered@chapter}[2][]{%
  \cleardoublepage
  \if@firstnumbered
    \pagenumbering{arabic}%
    \renewcommand{\thepage}{\arabic{page}}%
    \setcounter{page}{1}%
    \addtocontents{toc}{\protect\tocdivisionentry{Main Text}}%
    \addtocontents{toc}{\protect\setcounter{tocdepth}{1}}%
    \global\@firstnumberedfalse
  \fi
  \refstepcounter{chapter}%
  \setcounter{section}{0}%
  \if\relax\detokenize{#1}\relax
    \old@chapter*{{\chapterformat}#2}%
    \cleanmark{#2}%
    \addcontentsline{toc}{chapter}{\protect\numberline{\thechapter}#2}%
  \else
    \old@chapter*[#1]{{\chapterformat}#2}%
    \cleanmark{#1}%
    \addcontentsline{toc}{chapter}{\protect\numberline{\thechapter}#1}%
  \fi
}

% Definitive Chapter Entry layout in TOC (MLSysBook layout)
\renewcommand*{\l@chapter}[2]{%
  \ifnum \c@tocdepth >\m@ne
    \addpenalty{-\@highpenalty}%
    \vskip 0.60em\relax
    {%
      \leftskip 0pt\relax
      \rightskip \@tocrmarg
      \parfillskip -\rightskip
      \parindent 0pt\relax
      \@afterindenttrue
      \interlinepenalty\@M
      \leavevmode
      \begingroup
        \renewcommand{\numberline}[1]{{\normalfont\rmfamily\bfseries\fontsize{9.5pt}{12pt}\selectfont\color{ink}\format@toc@number{##1}\hspace{0.55em}}}%
        {\normalfont\rmfamily\bfseries\fontsize{9.5pt}{12pt}\selectfont\color{ink}#1}\nobreak
      \endgroup
      \TOCLineLeaderFill
      \nobreak
      {\normalfont\rmfamily\bfseries\fontsize{9.5pt}{12pt}\selectfont\color{ink}#2}%
      \par
    }%
  \fi
}

% Definitive Section Entry layout in TOC (MLSysBook layout)
\renewcommand*{\l@section}[2]{%
  \ifnum \c@tocdepth >\z@
    \vskip 0.10em\relax
    {%
      \leftskip 1.4em\relax
      \rightskip \@tocrmarg
      \parfillskip -\rightskip
      \parindent 0pt\relax
      \@afterindenttrue
      \interlinepenalty\@M
      \leavevmode
      \begingroup
        \renewcommand{\numberline}[1]{\makebox[2.9em][l]{##1\hspace{0.4em}}}%
        {\normalfont\rmfamily\fontsize{8.8pt}{11pt}\selectfont\color{ink}#1}\nobreak
      \endgroup
      \TOCLineLeaderFill
      \nobreak
      {\normalfont\rmfamily\fontsize{8.8pt}{11pt}\selectfont\color{ink}#2}%
      \par
    }%
  \fi
}
\makeatother


% ---------------------------------------------------------------------------
% 5. Floats: captions and tables
%
% Captions lead with a bold takeaway, so the label sits in structural blue and
% the caption text drops to secondary ink.
% ---------------------------------------------------------------------------

\setkomafont{captionlabel}{\normalfont\sffamily\bfseries\fontsize{8.2pt}{10.5pt}\selectfont\color{structural}}
\setkomafont{caption}{\normalfont\sffamily\fontsize{8.0pt}{10.2pt}\selectfont\color{softink}}

% Tables: Clean, academic, compact styling (Cambridge/MIT Press monograph style)
\renewcommand{\arraystretch}{1.18}
\setlength{\tabcolsep}{3.8pt}

% Table rules styling
\setlength{\heavyrulewidth}{1.0pt}
\setlength{\lightrulewidth}{0.5pt}

% Prevent awkward orphan headers: require at least 6 baselines of vertical space
% before starting a longtable, allowing tables to render directly under their introducing prose
\BeforeBeginEnvironment{longtable}{%
  \par\penalty0\needspace{6\baselineskip}%
}
\AtBeginEnvironment{longtable}{%
  \normalfont\sffamily\fontsize{7.6pt}{9.6pt}\selectfont%
}
\AfterEndEnvironment{longtable}{%
  \normalfont\selectfont%
}
\AtBeginEnvironment{tabular}{%
  \normalfont\sffamily\fontsize{7.6pt}{9.6pt}\selectfont%
  \setlength{\tabcolsep}{3.8pt}%
  \renewcommand{\arraystretch}{1.18}%
}
\AfterEndEnvironment{tabular}{%
  \normalfont\selectfont%
}
\AtBeginEnvironment{tabular*}{%
  \normalfont\sffamily\fontsize{7.6pt}{9.6pt}\selectfont%
  \setlength{\tabcolsep}{3.8pt}%
  \renewcommand{\arraystretch}{1.18}%
}
\AfterEndEnvironment{tabular*}{%
  \normalfont\selectfont%
}

% Floats: Academic textbook standard float placement with placeins barrier
\makeatletter
\def\fps@figure{tbp}
\def\fps@table{htbp}
\makeatother
\setcounter{topnumber}{3}
\setcounter{bottomnumber}{2}
\setcounter{totalnumber}{5}
\setcounter{dbltopnumber}{3}
\renewcommand{\topfraction}{0.90}       % max fraction of page for floats at top
\renewcommand{\bottomfraction}{0.60}    % allow bottom floats
\renewcommand{\textfraction}{0.10}      % min fraction of page for text
\renewcommand{\floatpagefraction}{0.65} % min fraction of float page occupied
\renewcommand{\dbltopfraction}{0.90}
\renewcommand{\dblfloatpagefraction}{0.65}
\setlength{\floatsep}{10pt plus 2pt minus 2pt}
\setlength{\textfloatsep}{12pt plus 2pt minus 4pt}
\setlength{\intextsep}{10pt plus 2pt minus 2pt}
\setlength{\dblfloatsep}{10pt plus 2pt minus 2pt}
\setlength{\dbltextfloatsep}{12pt plus 2pt minus 4pt}






% ---------------------------------------------------------------------------
% 6. Code listings
%
% commandchars is load bearing. Pandoc emits \CommentTok{...} and
% \NormalTok{...} inside Highlighting, and without commandchars those macros
% are typeset as literal text instead of applied. Redefining Highlighting
% without it silently breaks every syntax-highlighted block in the book,
% including the dossier YAML, which is exactly what happened before.
% ---------------------------------------------------------------------------

% Shaded environment: tighten padding and frame for listings
\ifcsname Shaded\endcsname\else
  \newenvironment{Shaded}{}{}
\fi
\renewenvironment{Shaded}{%
  \begin{tcolorbox}[
    enhanced,
    breakable,
    colback=paleink!10!white,
    colframe=paleink!30!white,
    boxrule=0.4pt,
    arc=1.0pt,
    left=3pt,
    right=3pt,
    top=2pt,
    bottom=2pt,
    before skip=3pt,
    after skip=4pt
  ]%
}{%
  \end{tcolorbox}%
}

\DefineVerbatimEnvironment{Highlighting}{Verbatim}{
  commandchars=\\\{\},
  fontsize=\fontsize{6.6pt}{8.0pt}\ttfamily\selectfont,
  breaklines=true,
  breakanywhere=true,
  xleftmargin=1pt,
  xrightmargin=1pt
}

% Plain verbatim carries no tokens, so it takes no commandchars.
\RecustomVerbatimEnvironment{verbatim}{Verbatim}{
  fontsize=\fontsize{6.6pt}{8.0pt}\ttfamily\selectfont,
  breaklines=true,
  breakanywhere=true,
  xleftmargin=1pt,
  xrightmargin=1pt
}


% ---------------------------------------------------------------------------
% 7. Margin notes and footnotes
%
% Footnotes are set in the margin (reference-location: margin in _quarto.yml),
% so both read as one register: small, sans, secondary ink, sized compactly
% for the narrow margin column.
% ---------------------------------------------------------------------------

\renewcommand*{\marginfont}{%
  \RaggedRight\normalfont\sffamily\fontsize{7.5pt}{9.8pt}\selectfont\color{softink}}
\setkomafont{footnote}{%
  \normalfont\sffamily\fontsize{7.8pt}{10.2pt}\selectfont\color{softink}}
\setkomafont{footnotemark}{%
  \normalfont\sffamily\bfseries\fontsize{7.8pt}{10.2pt}\selectfont\color{consequence}}

% Sidenotes package hook: Quarto emits \sidenote{\footnotesize ...} when
% reference-location: margin is used. We define custom handlers to neutralize
% pandoc's \footnotesize and enforce compact sans-serif typography in margin notes.
\makeatletter
\NewDocumentCommand \customsidenotetext { o o +m }
{%
  \IfNoValueOrEmptyTF{#1}%
    {%
      \marginnote{%
        \RaggedRight\normalfont\sffamily\fontsize{6.8pt}{8.8pt}\selectfont\color{softink}%
        {\bfseries\color{consequence}\textsuperscript{\thesidenote}}\nobreakspace
        {\let\footnotesize\relax #3}\par
      }%
      \refstepcounter{sidenote}%
    }%
    {%
      \marginnote{%
        \RaggedRight\normalfont\sffamily\fontsize{6.8pt}{8.8pt}\selectfont\color{softink}%
        {\bfseries\color{consequence}\textsuperscript{#1}}\nobreakspace
        {\let\footnotesize\relax #3}\par
      }%
    }%
}

\NewDocumentCommand \customsidenotemark { m }
{%
  \leavevmode
  \ifhmode
      \edef \@x@sf {\the \spacefactor }%
      \nobreak
  \fi
  \hbox {\textsuperscript{\normalfont\sffamily\bfseries\fontsize{6.8pt}{8.8pt}\selectfont\color{consequence}#1}}%
  \ifhmode
      \spacefactor \@x@sf
  \fi
  \relax
}

\AtBeginDocument{%
  \renewcommand*{\@sidenotes@placemarginal}[2]{%
    \IfNoValueOrEmptyTF{#1}{\marginnote{#2}}{\marginnote[#1]{#2}}%
  }%
  \let\sidenotetext\customsidenotetext
  \let\@sidenotes@thesidenotemark\customsidenotemark
}
\makeatother


% ---------------------------------------------------------------------------
% 8. Links
%
% Deferred to \AtBeginDocument so hyperref has loaded and Quarto has finished
% its own \hypersetup.
% ---------------------------------------------------------------------------

\AtBeginDocument{%
  \hypersetup{
    colorlinks=true,
    linkcolor=ethblue,
    urlcolor=ethpetrol,
    citecolor=ethpetrol,
    filecolor=ethpetrol
  }%
}


% ---------------------------------------------------------------------------
% 9. Utilities & Margin Definition Badges
% ---------------------------------------------------------------------------

% Margin Definition Badge: Sleek, compact card for margin glossary terms
\NewDocumentCommand{\defmargin}{ m +m }{%
  \marginnote{%
    \begin{tcolorbox}[
      enhanced,
      breakable=false,
      colback=structural!4!white,
      colframe=structural,
      boxrule=0pt,
      leftrule=2pt,
      arc=1pt,
      left=3.5pt,
      right=3pt,
      top=2pt,
      bottom=2.5pt,
      before skip=2pt,
      after skip=3pt
    ]%
      {\fontsize{6.5pt}{8pt}\sffamily\bfseries\scshape\color{structural}Definition}\\[1pt]
      {\fontsize{7.2pt}{9.2pt}\sffamily\bfseries\color{ink}#1}\\[1.5pt]
      {\fontsize{6.8pt}{8.8pt}\sffamily\color{softink}#2}%
    \end{tcolorbox}%
  }%
}

% fullwidth spans the text block plus the margin column, for an artifact the
% reader scans across rather than reads down.
%
%   \begin{fullwidth} ... \end{fullwidth}
%
% Use it sparingly, and only for genuinely wide material. A table of sentences
% does not become legible by getting wider; it needs shorter cells.
\newlength{\fullwidthlen}
\setlength{\fullwidthlen}{\dimexpr\textwidth+\marginparsep+\marginparwidth\relax}
\NewEnviron{fullwidth}{%
  \par\noindent
  \makebox[\textwidth][l]{%
    \begin{minipage}{\fullwidthlen}\BODY\end{minipage}%
  }\par
}


% ---------------------------------------------------------------------------
% 10. Margin Alignment Guides (development aid — toggle)
%
% Draws light frames around the text block, header/footer, and the margin-note
% column, so text, figures, tables, and notes can be aligned during authoring.
% Set \MarginDebugfalse for production builds; \MarginDebugtrue to show guides.
% ---------------------------------------------------------------------------

\newif\ifMarginDebug
\MarginDebugfalse% <<< set to \MarginDebugtrue to show margin alignment guides
\ifMarginDebug
  \usepackage{showframe}
  \renewcommand*\ShowFrameLinethickness{0.2pt}
  \renewcommand*\ShowFrameColor{\color{red!45}}
\fi



